% Socle Identitaire V14-FR — Version académique française
% Version originale en anglais : identity_bedrockV14.tex
% Sans PBAT, sans présence physique.
% Compilation : pdflatex -> pdflatex

\documentclass[11pt,a4paper]{article}

\usepackage[utf8]{inputenc}
\usepackage[T1]{fontenc}
\usepackage{lmodern}
\usepackage{microtype}
\usepackage[french]{babel}
\usepackage{amsmath,amssymb,amsthm}
\usepackage{mathtools}
\usepackage{booktabs}
\usepackage{array}
\usepackage{enumitem}
\usepackage[colorlinks=true,linkcolor=black,citecolor=black,urlcolor=blue]{hyperref}
\usepackage{geometry}

\geometry{a4paper,margin=25mm}

\newtheorem{theoreme}{Théorème}[section]
\newtheorem{definition}[theoreme]{Définition}
\newtheorem{remarque}[theoreme]{Remarque}
\newtheorem{corollaire}[theoreme]{Corollaire}
\renewcommand{\proofname}{Démonstration}

\title{\textbf{Socle Identitaire : Opérationnalisation du Substrat\\
Cryptographique Minimal de Rompel\\
pour l'Identité Souveraine}}
\author{Jean-François Bouchard\\
  \small PrivateNess Network (co-fondateur, DT)\\
  \small Emercoin (contributeur communautaire)\\
  \small \texttt{[contact via privateness.network]}}
\date{Mars 2026}

\begin{document}

\maketitle

\noindent\textbf{Divulgation relative à l'IA.} L'auteur a eu recours à des
outils de traitement automatique du langage pour la rédaction et la révision
du manuscrit.  L'ensemble des contributions techniques, la conception du
système, les réalisations, le cadrage théorique et les idées directrices
décrits ci-après sont issus de l'auteur, qui s'appuie sur neuf années
d'expérience opérationnelle dans le déploiement d'infrastructures fondées sur
Emercoin.

\begin{abstract}
Le Socle Identitaire n'a pas émergé de la théorie.  Il a émergé d'une
exigence~: éliminer toute clé institutionnelle en dessous de l'individu.  En
suivant cette exigence jusqu'à son plancher logique---en nous demandant quel
est le primitif cryptographique minimal permettant l'authentification
publique---nous
avons abouti aux fonctions à sens unique comme fondement irréductible.  La
démonstration de Rompel, en 1990, que les FSU sont nécessaires et suffisantes
pour les schémas de signature à infalsifiabilité existentielle sous attaque
adaptative à messages choisis~\cite{rompel1990} a été découverte \emph{après}
cette convergence.  Il avait déjà prouvé que nous nous tenions sur le socle.

Le système découple la cryptographie d'identité du consensus du répertoire
au moyen d'une couche d'ancrage en ajout exclusif, indépendante de la chaîne
de blocs, permettant la portabilité entre substrats hétérogènes tout en
préservant la vérifiabilité.  Nous instancions cette construction à l'aide du Stockage
Nom--Valeur (SNV) d'Emercoin ancré à Bitcoin par Preuve de Travail Auxiliaire
(AuxPoW), avec un déploiement publiquement vérifiable opérationnel depuis 2020.

Nous montrons en outre que des capacités système supplémentaires---nommage,
représentation structurée, récupération par contenu, exécution déterministe,
connectivité en réseau superposé, interaction économique---émergent comme
propriétés dérivées de ce substrat identitaire, sans nécessiter de nouvelles
racines de confiance.

La contribution n'est pas un nouveau primitif, mais une composition de
primitifs existants en un système qui réduit les hypothèses de confiance au
minimum irréductible, tout en restant déployable sur des infrastructures
matérielles grand public.
\end{abstract}

\noindent\textbf{Mots-clés :} fonctions à sens unique $\cdot$ signatures
numériques $\cdot$ identité souveraine $\cdot$ identité sur chaîne de blocs
$\cdot$ architecture agnostique $\cdot$ annuaire en ajout exclusif $\cdot$
migration post-quantique

%--------------------------------------------------------------------
\section{Introduction}
%--------------------------------------------------------------------

\subsection{Le problème : la capture institutionnelle}

Tout système d'identité numérique déployé à ce jour place des clés
institutionnelles en dessous de l'utilisateur~: autorités de certification,
fournisseurs d'identité, registres de clés, opérateurs de plateforme.  Ces
entités peuvent révoquer, refuser, surveiller ou contrôler l'identité sans le
consentement de l'utilisateur.

Le modèle de menace qui motive ce travail suppose un adversaire disposant de
ressources conventionnelles substantielles, d'un soutien politique potentiel
et d'un contrôle total sur l'infrastructure d'identité traditionnelle.  La
question centrale est la suivante~: \emph{l'identité peut-elle exister sans
clé institutionnelle en dessous de l'individu~?}

\subsection{La réponse théorique : Rompel 1990}

Rompel a démontré que les FSU sont à la fois nécessaires et suffisantes pour
les schémas $\mathrm{SIG}_{\mathrm{EUF\text{-}AAMC}}$~\cite{rompel1990}.
Il s'agit du minimum prouvable~: aucune hypothèse strictement plus faible ne
suffit pour l'authentification publique sous les hypothèses standard.

Notre thèse est la suivante~: en fixant la FSU comme plancher cryptographique
et en n'ajoutant qu'un annuaire en ajout exclusif, indépendant de toute chaîne
de blocs particulière, on obtient une construction identitaire qui réduit à
zéro la capture institutionnelle sous la clé individuelle.

\subsection{L'idée directrice}

Ce travail s'inscrit dans un motif récurrent de l'innovation cryptographique~:
la valeur ne réside pas dans l'invention de nouveaux primitifs, mais dans la
reconnaissance que des primitifs existants, liés d'une manière spécifique,
produisent des propriétés émergentes qu'aucun d'eux ne présente isolément.
Chaque composant du Socle Identitaire---signatures fondées sur les FSU,
annuaires en ajout exclusif---existe indépendamment dans la littérature
antérieure.  La contribution réside dans la liaison et dans l'argumentation
formelle des propriétés de sécurité que cette liaison engendre.

\subsection{Contributions}

\begin{enumerate}[label=\textbf{C\arabic*.},leftmargin=*]
  \item \textbf{Théorie.}  Formalisation d'un système d'identité indépendant
    de toute chaîne de blocs particulière, ancré dans les FSU à la manière de
    Rompel, avec
    des définitions en ligne et un théorème de solidité de l'authentification
    démontré par réduction de sécurité explicite.

  \item \textbf{Systèmes.}  Architecture~: couche d'identité fondée sur les FSU
    (Ed25519 par défaut~; cible de migration ML-DSA-65~\cite{fips204}) avec
    substrat de répertoire enfichable.  Mise en \oe{}uvre~: SNV d'Emercoin
    ancré à Bitcoin via AuxPoW~\cite{emercoin} et NESS Fiber~\cite{onchain}.
    Déployé et publiquement vérifiable depuis 2020.

  \item \textbf{Émergence systémique.}  Démonstration que le nommage, la
    représentation structurée, l'adressage par contenu, l'exécution
    déterministe, la connectivité en réseau superposé et l'interaction
    économique émergent tous comme propriétés dérivées du substrat identitaire,
    sans racines de confiance supplémentaires.

  \item \textbf{Analyse du coût d'attaque.}  Quantification du coût d'attaque
    sur les composants cryptographique et de répertoire à l'aide du cadre
    hack51~\cite{hack51}.
\end{enumerate}

\subsection{Plan de l'article}

La section~2 présente les préliminaires.  La section~3 établit le modèle de
menace.  La section~4 développe le modèle formel.  La section~5 décrit
l'architecture.  La section~6 démontre la solidité de l'authentification.
La section~7 présente les propriétés système émergentes.  La section~8
traite de l'instanciation concrète et de l'analyse du coût d'attaque.  La
section~9 décrit le système déployé.  La section~10 examine les travaux
connexes.  La section~11 discute les limites et perspectives.

%--------------------------------------------------------------------
\section{Préliminaires}
%--------------------------------------------------------------------

\subsection{Fonctions à sens unique et signatures numériques}

Une fonction $f : \{0,1\}^* \to \{0,1\}^*$ est une
\emph{fonction à sens unique} (FSU) si elle est calculable en temps
polynomial et si, pour tout adversaire probabiliste en temps polynomial (PPT)
$\mathcal{A}$,
\[
  \Pr\bigl[f(x') = f(x) \;\big|\; x \gets \{0,1\}^\lambda,\;
    x' \gets \mathcal{A}(1^\lambda, f(x))\bigr] \leq \mathrm{negl}(\lambda),
\]
où $\mathrm{negl}(\cdot)$ désigne une fonction négligeable.

Rompel~\cite{rompel1990}, complétant le programme de
Naor--Yung~\cite{naoryung1989} et ultérieurement de Haitner et
al.~\cite{haitner2010}, a démontré que les FSU sont \emph{nécessaires et
suffisantes} pour les schémas de signature à sécurité EUF-AAMC :
\[
  \mathrm{FSU} \iff \mathrm{SIG}_{\mathrm{EUF\text{-}AAMC}}.
\]

Cette équivalence signifie que l'existence d'un schéma EUF-AAMC-sécurisé
contraint le système à exactement une hypothèse d'intractabilité---l'existence
des FSU---et qu'aucune hypothèse plus faible n'est suffisante.  Le Socle
Identitaire opérationnalise ce résultat comme principe de conception~: la
paire de clés de l'individu est l'unique racine cryptographique de l'identité.

\subsection{Architectures d'identité traditionnelles}

L'infrastructure à clé publique (ICP) / X.509~\cite{rfc5280} place une
hiérarchie d'autorités de certification en dessous de la clé individuelle,
introduisant des oracles de révocation et des compromissions d'AC comme modes
de défaillance supplémentaires.  Les systèmes d'identité fédérés (OAuth~2.0,
OpenID Connect) délèguent l'identité à un fournisseur qui peut désauthentifier
les utilisateurs unilatéralement.  Même dans les contextes de chaînes de
blocs, les multisignatures de gouvernance et les clés administratives se
trouvent souvent en dessous de l'individu dans la pile de confiance.  Ces
choix architecturaux ajoutent des hypothèses au-delà des FSU sans accroître
la sécurité cryptographique de l'authentification elle-même.

\subsection{AuxPoW Bitcoin et minage fusionné}

La Preuve de Travail Auxiliaire permet le minage fusionné d'Emercoin avec
Bitcoin~: Emercoin hérite ainsi d'une fraction de la sécurité de Bitcoin pour
des applications non monétaires.  Khovayko~\cite{hack51} introduit le
\emph{coefficient hack51}---défini comme $c_{\mathrm{att}} / c_{\mathrm{cap}}$---comme métrique numériquement mesurable de
résistance à l'attaque à 51\,\%.  Sous cette métrique, le triple consensus
d'Emercoin (PoS + PoW + minage fusionné) produit un coefficient nettement
supérieur à 1,0~: attaquer Emercoin exige des ressources approchant cinq fois
sa capitalisation boursière.  Nous appliquons ce cadre à la section~\ref{sec:instanciation}.

%--------------------------------------------------------------------
\section{Modèle de menace}
%--------------------------------------------------------------------

\subsection{Capacités de l'adversaire}

\textbf{Hypothèses retenues.}  Ressources conventionnelles substantielles et
soutien politique~; contrôle total de l'infrastructure d'identité traditionnelle
(AC, fournisseurs d'identité, registres KYC)~; capacité à interdire des
logiciels, saisir des biens et contraindre des individus~; adversaire
calculatoire classique (en temps polynomial).

\textbf{Hypothèses non retenues.}  Ne peut pas violer les FSU sous les
hypothèses de complexité standard~; ne peut pas mener une attaque à 51\,\%
soutenue sur Bitcoin (quantifiée à la section~\ref{sec:instanciation})~; ne
peut pas compromettre simultanément une infrastructure distribuée à l'échelle
mondiale.

\textbf{Périmètre explicite.}  Ce travail traite des adversaires classiques.
La migration post-quantique vers ML-DSA (FIPS~204~\cite{fips204}) est abordée
à la section~\ref{sec:discussion}.

\subsection{Exigences de conception}

\begin{description}[leftmargin=*,style=nextline]
  \item[E1 --- Irréductibilité mathématique.]  FSU comme plancher
    cryptographique.  Aucun primitif strictement plus faible en dessous.
  \item[E2 --- Absence de clé institutionnelle.]  La clé individuelle est
    l'unique racine cryptographique de confiance.  Aucune AC, fournisseur
    d'identité ou registre en dessous de l'utilisateur.
  \item[E3 --- Agnosticisme vis-à-vis de la chaîne.]  Cryptographie d'identité
    indépendante de tout consensus de répertoire particulier.  Portable entre
    substrats en ajout exclusif.
  \item[E4 --- Sécurité quantifiable.]  Coûts d'attaque explicites et
    mesurables.
  \item[E5 --- Vérification opérationnelle.]  Système déployé, non simple
    proposition théorique.  Vérification publique de toutes les affirmations.
\end{description}

%--------------------------------------------------------------------
\section{Modèle formel}
%--------------------------------------------------------------------

Soit $\lambda \in \mathbb{N}$ le paramètre de sécurité et
$\mathrm{negl}(\lambda)$ une fonction négligeable.

\begin{definition}[Schéma de signature]
  Un \emph{schéma de signature} est un triplet $\Sigma = (\mathsf{KeyGen},
  \mathsf{Sign}, \mathsf{Verify})$ d'algorithmes PPT où :
  \begin{itemize}
    \item $\mathsf{KeyGen}(1^\lambda) \to (sk, pk)$,
    \item $\mathsf{Sign}(sk, m) \to \sigma$ pour $m \in \{0,1\}^*$,
    \item $\mathsf{Verify}(pk, m, \sigma) \to \{0,1\}$.
  \end{itemize}
  $\Sigma$ est \emph{correct} si $\mathsf{Verify}(pk, m, \mathsf{Sign}(sk,m))
  = 1$ pour tout $(sk,pk) \gets \mathsf{KeyGen}(1^\lambda)$ et tout $m$.
\end{definition}

\begin{definition}[Sécurité EUF-AAMC]
  $\Sigma$ est à \emph{infalsifiabilité existentielle sous attaque adaptative
  à messages choisis} si pour tout adversaire PPT $\mathcal{A}$,
  \[
    \Pr\bigl[\mathsf{EUF\text{-}AAMC}(\mathcal{A},\Sigma,\lambda) = 1\bigr]
    \leq \mathrm{negl}(\lambda),
  \]
  où $\mathsf{EUF\text{-}AAMC}(\mathcal{A},\Sigma,\lambda)$ est le jeu
  standard~: $(sk,pk) \gets \mathsf{KeyGen}(1^\lambda)$~;
  $\mathcal{A}^{\mathsf{Sign}(sk,\cdot)}(pk)$ produit $(m^*, \sigma^*)$~;
  $\mathcal{A}$ gagne si $\mathsf{Verify}(pk, m^*, \sigma^*) = 1$
  et $m^*$ n'a jamais été soumis à l'oracle de signature.
\end{definition}

\begin{definition}[Annuaire public]
  Un \emph{annuaire public} $\mathcal{D}$ est un ensemble de paires
  $(id, pk) \in \mathcal{ID} \times \mathcal{CP}$ muni du prédicat
  d'appartenance
  \[
    \mathsf{Lier}(id, pk) = 1 \iff (id, pk) \in \mathcal{D}.
  \]
  $\mathcal{D}$ est \emph{en ajout exclusif} si, pour tout $(id, pk)$
  une fois inséré, il n'existe aucune procédure efficace pouvant rendre
  $\mathsf{Lier}(id, pk)$ égal à $0$.
\end{definition}

\begin{definition}[$\varepsilon_{\mathrm{dir}}$-Finalité]
  Un annuaire en ajout exclusif $\mathcal{D}$ possède la
  \emph{$\varepsilon_{\mathrm{dir}}$-finalité} si, pour tout adversaire PPT
  $\mathcal{A}$ et toute entrée $(id, pk)$ inscrite dans $\mathcal{D}$ à la
  profondeur $d$,
  \[
    \Pr\bigl[\mathcal{A}\ \text{provoque } \mathsf{Lier}(id,pk) = 0\bigr]
    \leq \varepsilon_{\mathrm{dir}}(\lambda, d).
  \]
  Lorsque la profondeur de confirmation $d \to \infty$,
  $\varepsilon_{\mathrm{dir}} \to 0$.
\end{definition}

\begin{definition}[Prédicat d'authentification]
  Étant donné l'annuaire $\mathcal{D}$, le \emph{prédicat d'authentification}
  est
  \[
    \mathsf{Auth}(id, c, \sigma) := 1 \iff
    \exists\, pk :\;
    \mathsf{Lier}(id, pk) = 1
    \;\wedge\;
    \mathsf{Verify}(pk, c, \sigma) = 1.
  \]
\end{definition}

%--------------------------------------------------------------------
\section{Architecture}
%--------------------------------------------------------------------

\subsection{Couche~1 : Cryptographie d'identité}

Un participant génère $(sk, pk) \gets \mathsf{KeyGen}(1^\lambda)$ à partir
d'une graine de haute entropie.  La clé publique $pk$, associée à un
identifiant choisi $id \in \mathcal{ID}$, constitue l'objet identitaire
$\mathcal{I} = (id, pk)$.  L'authentification auprès de tout vérificateur
capable d'évaluer $\mathsf{Lier}$ se réduit alors à un protocole
défi--réponse standard~:

\begin{enumerate}[leftmargin=*]
  \item Le vérificateur tire $c \gets \{0,1\}^\lambda$ et l'envoie au prouveur.
  \item Le prouveur calcule $\sigma \gets \mathsf{Sign}(sk, c)$ et renvoie $\sigma$.
  \item Le vérificateur accepte si et seulement si $\mathsf{Auth}(id, c, \sigma) = 1$.
\end{enumerate}

Aucun certificat, jeton de session ou oracle tiers n'intervient dans cet échange.

\subsection{Couche~2 : Annuaire en ajout exclusif}

L'annuaire $\mathcal{D}$ lie les identifiants aux clés publiques.  Les
exigences sont exactement celles de la Définition~3 : vérifiabilité globale,
persistance et $\varepsilon_{\mathrm{dir}}$-finalité.  La couche d'identité
et la couche de répertoire sont découplées~: tout substrat satisfaisant la
Définition~3 constitue une instanciation valide.  La section~\ref{sec:instanciation}
en fournit une concrète.

%--------------------------------------------------------------------
\section{Analyse de sécurité}
%--------------------------------------------------------------------

\begin{theoreme}[Solidité de l'authentification]
  \label{thm:solidite}
  Soit $\Sigma$ à sécurité EUF-AAMC avec borne d'avantage
  $\varepsilon_{\mathrm{sig}}$, et soit $\mathcal{D}$ un annuaire en ajout
  exclusif avec $\varepsilon_{\mathrm{dir}}$-finalité.  Pour tout adversaire
  PPT $\mathcal{A}$ ne détenant pas $sk$, la probabilité que $\mathcal{A}$
  amène un vérificateur à accepter sous l'identité $id$ est au plus
  $\varepsilon_{\mathrm{sig}}(\lambda) + \varepsilon_{\mathrm{dir}}(\lambda, d)$.
\end{theoreme}

\begin{proof}
Supposons que l'adversaire $\mathcal{A}$ gagne le jeu d'authentification~:
il produit un triplet valide $(id, c, \sigma)$ sans détenir $sk$.

\textbf{Réduction.}  Construisons un faussaire $\mathcal{B}$ contre $\Sigma$.
$\mathcal{B}$ reçoit la clé publique de défi $\hat{pk}$.  Il simule l'annuaire
pour $\mathcal{A}$ en répondant aux requêtes $\mathsf{Lier}$ portant sur $id$
par $\hat{pk}$.\footnote{$\mathcal{B}$ simule l'annuaire localement~; il
n'écrit rien sur un registre réel.}  Nous distinguons deux cas mutuellement
exclusifs.

\textbf{Cas~1 : l'intégrité de l'annuaire est préservée} (probabilité
$\geq 1 - \varepsilon_{\mathrm{dir}}$).  La seule clé $pk$ satisfaisant
$\mathsf{Lier}(id,pk)=1$ est $\hat{pk}$.  Lorsque $\mathcal{A}$ déclenche
l'authentification, $\mathcal{B}$ lui transmet le défi EUF-AAMC $c$.  Étant
donné que $c$ est tiré fraîchement après que $\mathcal{A}$ a fixé sa stratégie
et n'a jamais été soumis à l'oracle $\mathsf{Sign}(sk,\cdot)$ légitime,
tout $\sigma$ valide renvoyé par $\mathcal{A}$ est une falsification
existentielle.  Par sécurité EUF-AAMC de $\Sigma$ :
\[
  \Pr[\mathcal{A}\ \text{gagne} \mid \text{Cas~1}]
  \leq \varepsilon_{\mathrm{sig}}(\lambda).
\]

\textbf{Cas~2 : l'intégrité de l'annuaire est compromise} (probabilité
$\leq \varepsilon_{\mathrm{dir}}$).  $\mathcal{A}$ remplace $(id, \hat{pk})$
par $(id, pk')$ pour une clé dont il contrôle la clé secrète correspondante.
Cet événement se produit avec probabilité au plus $\varepsilon_{\mathrm{dir}}(\lambda,d)$
d'après la Définition~4.  Les attaques sur l'annuaire ramènent le coût à la
couche de répertoire, quantifiée à la section~\ref{sec:instanciation}.

Par union :
$\Pr[\mathcal{A}\ \text{gagne}] \leq
\varepsilon_{\mathrm{sig}} + \varepsilon_{\mathrm{dir}}$,
toutes deux négligeables sous les hypothèses énoncées.
\end{proof}

\begin{corollaire}
  Toute usurpation réussie de l'identité $id$ exige soit (i) de violer $\Sigma$
  avec avantage $> \varepsilon_{\mathrm{sig}}$, soit (ii) de corrompre
  l'annuaire à un coût excédant le seuil économique impliqué par
  $\varepsilon_{\mathrm{dir}}$.  Aucun secret institutionnel n'apparaît dans
  l'une ou l'autre condition.
\end{corollaire}

%--------------------------------------------------------------------
\section{Propriétés système émergentes}
\label{sec:derivees}
%--------------------------------------------------------------------

Dès lors qu'un substrat identitaire satisfaisant les Définitions~1--5 existe,
des capacités système supplémentaires deviennent disponibles comme propriétés
dérivées, sans introduction de nouvelles racines de confiance.

\subsection{Nommage et résolution}

Les liaisons en ajout exclusif entre noms intelligibles et clés publiques
mettent en \oe{}uvre un nommage décentralisé.  La résolution est une
vérification d'identité~: pour résoudre le nom $n$, récupérer $(n, pk)$
dans $\mathcal{D}$ et vérifier $\mathsf{Lier}(n, pk) = 1$.  Aucune autorité
hiérarchique n'est requise.  Cette approche subsume la fonctionnalité des
propositions de nommage sur chaîne de blocs~\cite{handshake} sans modèle de
confiance séparé.

\subsection{Représentation structurée des enregistrements}

Un schéma JSON canonique en ajout exclusif---le format WORM (\emph{Write-Once
Read-Many})~\cite{worm}---définit des enregistrements typés stockés comme
valeurs dans l'annuaire.  La sérialisation déterministe permet à des n\oe{}uds
indépendants d'analyser, d'indexer et d'interroger les enregistrements sans
autorité de schéma partagée.

\subsection{Adressage par contenu}

La récupération de contenu utilise des hachés cryptographiques~: un n\oe{}ud
stocke une donnée $m$ à l'adresse $h = H(m)$ et tout récupérateur vérifie
$H(m') = h$ avant utilisation.  L'intégrité est indépendante du
transport~\cite{ipfs}.  Combinée à la couche d'identité, une partie
publiante peut signer une adresse de contenu, liant l'intégrité de la
récupération à l'authenticité de l'identité.

\subsection{Exécution déterministe}

Un calcul est modélisé comme une fonction pure $F$ appliquée à des entrées
adressées par contenu~: étant données les entrées $\{h_i\}$ et un programme
adressé par contenu $h_F$, la sortie $h_{\mathrm{sortie}} = H(F(m_1,
\ldots, m_k))$ est déterministe et indépendamment reproductible par tout
n\oe{}ud capable de récupérer les entrées.  La vérification de l'exécution
se réduit à vérifier que le haché de sortie publié correspond à l'évaluation
locale indépendante.

\subsection{Connectivité en réseau superposé}

La découverte de pairs et l'acheminement de messages peuvent être mis en
\oe{}uvre sur des points d'extrémité adressés par identité.  Chaque point
d'extrémité étant authentifié via $\mathsf{Auth}$, le réseau superposé
hérite de la sécurité de la couche d'identité et ne requiert pas d'ICP
séparée pour l'authentification du transport.

\subsection{Interaction économique}

Des modèles de récompense et de paiement fondés sur la participation peuvent
être construits à l'aide de l'attribution liée à l'identité.  Chaque action
économique est signée par $sk$ de l'acteur, rendant les contributions
attribuables et auditables sans opérateur de registre centralisé.

%--------------------------------------------------------------------
\section{Instanciation concrète}
\label{sec:instanciation}
%--------------------------------------------------------------------

\subsection{Schéma de signature}

Nous instancions $\Sigma$ avec Ed25519~\cite{eddsa}, un schéma de type
Schnorr sur Curve25519 offrant une sécurité d'environ $2^{128}$ contre les
attaques connues.  La génération de clés utilise une graine multi-source à
haute entropie~: plusieurs blocs de 1024~octets issus de \texttt{os.urandom}
(1\,536~bits d'état interne total), la sortie RC4OK du générateur de nombres aléatoires (GNA) du c\oe{}ur Emercoin, la gigue temporelle du processeur et des identifiants matériels,
mélangés via UHEPRNG~\cite{uheprng}.

\subsection{Répertoire : SNV Emercoin + AuxPoW}

L'annuaire public $\mathcal{D}$ est mis en \oe{}uvre à l'aide du Stockage
Nom--Valeur (SNV) d'Emercoin~\cite{emercoin}, un magasin clé--valeur natif de la chaîne de blocs.  Les blocs Emercoin sont minés de façon fusionnée avec
Bitcoin via AuxPoW~\cite{hack51}, ancrant l'état de l'annuaire à la preuve
de travail cumulée de Bitcoin.

Les enregistrements d'identité sont écrits comme objets JSON WORM sous des
clés SNV de la forme \texttt{id:$\langle$identifiant$\rangle$}.  Une fois
écrits, les enregistrements peuvent être renouvelés mais non supprimés,
satisfaisant l'exigence d'ajout exclusif de la Définition~3.

\subsection{Analyse du coût d'attaque}

\begin{center}
\begin{tabular}{lll}
  \toprule
  \textbf{Vecteur d'attaque} & \textbf{Condition} & \textbf{Résultat} \\
  \midrule
  Violer Ed25519       & Violer l'hypothèse FSU & Infaisable ($2^{128}$ op.) \\
  Réécrire le SNV Emercoin & Maj. PoS + 51\,\% AuxPoW & 5$\times$ cap. EMC~\cite{hack51} \\
  Compromettre NESS Fiber  & Vol des clés de publication & Déni de service uniquement ; aucune falsification \\
  Contraindre un utilisateur & Coercition physique / légale & Hors du périmètre crypto \\
  \bottomrule
\end{tabular}
\end{center}

\textbf{Falsification de signature.}  Violer Ed25519 exige de résoudre le
problème du logarithme discret sur Curve25519~; aucun algorithme sous-exponentiel
n'est connu ($O(2^{128})$ opérations, infaisable classiquement).

\textbf{Réécriture de l'annuaire.}  Modifier une entrée SNV finalisée exige
une attaque à $>50\,\%$ sur la chaîne AuxPoW d'Emercoin.  Sous le cadre
hack51~\cite{hack51}, ce coût approche cinq fois la capitalisation boursière
d'Emercoin.  Nuzzi et al.~\cite{nuzzi2024} estiment le coût d'une attaque
soutenue contre Bitcoin à entre 5 et 20~milliards USD, fournissant une borne
inférieure sur la seule composante AuxPoW.

\textbf{Compromission des clés NESS Fiber.}  Une clé de publication volée
ne peut pas réécrire les blocs passés (la BoltDB en ajout exclusif rend
impossible la réécriture de l'historique).  La seule action disponible est
une réinitialisation vers le bloc de genèse, immédiatement visible pour tous
les participants.  La compromission de clé est donc immédiatement détectable dans tous les cas.

%--------------------------------------------------------------------
\section{Démonstration opérationnelle}
%--------------------------------------------------------------------

Le système est opérationnel depuis 2020 sur le réseau PrivateNess.
L'enregistrement d'identité canonique en chaîne~\cite{onchain}---adresse de
genèse, clé publique de chaîne de blocs et signature de genèse---est ancré
dans le bloc SNV Emercoin \#445\,669 (1\textsuperscript{er}~novembre~2020,
22:16:45~UTC), transitivement ancré à Bitcoin via AuxPoW, et
indépendamment auditable à l'adresse
\url{https://explorer.emercoin.com/dpo/PrivateNESS.Network} par tout tiers
sans recours à un intermédiaire de confiance.  Ceci n'est pas un article
décrivant ce qui pourrait être construit.  C'est un article décrivant ce qui
fonctionne depuis 2020.

\textbf{NESS Fiber.}  NESS Fiber emploie le consensus Fiber de type
Skycoin~\cite{skycoin} avec des clés de publication annoncées dans le SNV
Emercoin, créant une chaîne de confiance vérifiable~: la sécurité de NESS
Fiber s'ancre dans le SNV Emercoin, qui s'ancre à Bitcoin via AuxPoW.  Tout
vérificateur peut confirmer indépendamment la liaison à l'explorateur de
blocs Emercoin, sans recours à un intermédiaire.

\textbf{Flux de résolution.}
\begin{enumerate}[leftmargin=*]
  \item Rechercher \texttt{id:$\langle$nom$\rangle$} dans le SNV Emercoin~;
    vérifier $\mathsf{Lier}(\mathit{nom}, pk) = 1$.
  \item Récupérer le contenu associé par haché~;
    vérifier $H(m') = h$.
  \item Vérifier la signature du publieur~:
    $\mathsf{Verify}(pk, h, \sigma) = 1$.
  \item Exécuter le programme adressé par contenu de façon déterministe.
\end{enumerate}

Chaque étape est indépendamment vérifiable par tout n\oe{}ud disposant d'un
accès à l'annuaire public et aux fonctions de hachage standard.  Aucun
intermédiaire de confiance n'y participe.  Le système fonctionne sur du
matériel grand public incluant des ordinateurs à carte unique de classe
RK3528 et des systèmes grand public réaffectés, avec un coût de déploiement
complet vérifié inférieur à 50~CAD~\cite{onchain}.

%--------------------------------------------------------------------
\section{Travaux connexes}
\label{sec:connexes}
%--------------------------------------------------------------------

\textbf{X.509 et ICP.}  La RFC~5280~\cite{rfc5280} place une hiérarchie
d'AC en dessous des clés individuelles.  Les AC peuvent émettre des
certificats erronés, être compromises ou contraintes, violant l'exigence
qu'aucune clé institutionnelle ne se trouve en dessous de
l'individu~\cite{goldreich2001}.  Le Socle Identitaire supprime entièrement
cette couche.

\textbf{Identifiants décentralisés (DID).}  La spécification W3C
DID~\cite{did2022} définit un cadre générique pour les identifiants
auto-souverains.  Le Socle Identitaire est compatible avec le modèle DID au
niveau de la couche d'identifiant, mais fournit un fondement formel de
sécurité plus robuste et un argument de finalité économique concret absent de
la spécification de base des DID.

\textbf{FIDO2 / WebAuthn.}  La norme FIDO2~\cite{webauthn2021} élimine les
mots de passe en utilisant des clés asymétriques liées aux appareils, s'alignant
étroitement sur le primitif d'identité de la Définition~1.  Cependant,
FIDO2 s'appuie sur des certificats d'attestation des fabricants de matériel,
réintroduisant une ancre de confiance institutionnelle.  Le Socle Identitaire
ne requiert aucune couche d'attestation.

\textbf{Connexion avec Ethereum (EIP-4361).}  L'EIP-4361~\cite{eip4361} lie
l'identité aux clés de comptes Ethereum et utilise l'état Ethereum comme
répertoire.  Cela couple la finalité de l'identité au consensus et à la
gouvernance Ethereum, et restreint le schéma de signature à secp256k1.  Le
Socle Identitaire généralise ce motif à des schémas EUF-AAMC arbitraires et
à des annuaires en ajout exclusif arbitraires.

\textbf{ENS et Handshake.}  L'Ethereum Name Service (ENS) et
Handshake~\cite{handshake} proposent un nommage décentralisé mais ne
définissent pas de substrat d'authentification formel.  Ils peuvent être vus
comme des instanciations concrètes de la couche de répertoire $\mathcal{D}$,
compatibles avec l'architecture présentée ici.

\textbf{Attestations vérifiables.}  Le modèle W3C d'attestations
vérifiables~\cite{vc2019} prend en charge les déclarations signées par un
émetteur, introduisant ainsi les clés d'émetteur comme ancres de confiance.
Le Socle Identitaire opère une couche en dessous des attestations vérifiables~:
il fournit la liaison d'identité à partir de laquelle la clé du détenteur
d'une attestation peut être dérivée.

%--------------------------------------------------------------------
\section{Discussion}
\label{sec:discussion}
%--------------------------------------------------------------------

\textbf{Ce que le socle fournit.}  Une liaison vérifiable et infalsifiable
entre un identifiant et une clé publique, sécurisée par des hypothèses
cryptographiques standard et un argument de finalité économique quantifié.
Toutes les propriétés de couche supérieure décrites à la
section~\ref{sec:derivees} héritent de cette base de sécurité.

\textbf{Ce qui est hors périmètre.}  Le modèle ne traite pas~:
\begin{itemize}[leftmargin=*]
  \item \emph{Compromission du point d'extrémité et logiciels malveillants.}
    Un adversaire ayant accès à l'appareil détenant $sk$ compromet
    l'authentification indépendamment du protocole.
  \item \emph{Ingénierie sociale et hameçonnage.}  Attaques hors bande
    incitant le détenteur légitime de la clé à signer des défis adversariaux.
  \item \emph{Coercition légale et réglementaire.}  La divulgation forcée de
    clé est orthogonale à la sécurité cryptographique.
\end{itemize}
Ces préoccupations sont opérationnelles et sociales ; elles requièrent des
couches supplémentaires et sont délibérément hors périmètre.

\textbf{Comparaison avec les approches à divulgation nulle.}  Le protocole
d'authentification de la section~5.1 révèle la paire défi--réponse
$(c, \sigma)$, ce qui peut permettre une liaison entre sessions si $pk$ est
public.  Les schémas d'accréditations anonymes~\cite{goldreich2001} peuvent
fournir une impossibilité de liaison au prix d'une complexité accrue~; il
s'agit d'une extension naturelle du socle plutôt que d'une contradiction.

\textbf{Migration post-quantique.}  Ed25519 est vulnérable à l'algorithme de
Shor.  La trajectoire de migration envisagée est ML-DSA-65 (CRYSTALS-Dilithium3,
FIPS~204~\cite{fips204})~:
\begin{itemize}[leftmargin=*]
  \item Remplacer Ed25519 par ML-DSA-65 (clé publique de 1\,952~octets~;
    prise en charge par le champ de valeur EmerNVS).
  \item Le schéma \texttt{<keys>} WORM~\cite{worm} prend en charge nativement
    plusieurs algorithmes~; la migration ajoute
    \texttt{<key algo="mldsa65">} sans modifier les enregistrements existants.
  \item L'ancrage dans l'annuaire, la structure NESS Fiber et le protocole
    d'authentification demeurent inchangés.
  \item OpenSSL~3.5 (avril~2025) intègre la prise en charge native de ML-DSA
    dans Armbian Linux en branche principale.
\end{itemize}
Les couches de répertoire et de connectivité n'exigent aucune migration cryptographique~;
seul le schéma de signature de la couche~1 requiert une migration.

\textbf{Rotation des clés.}  Le modèle de répertoire prend en charge la
rotation de clé en enregistrant une nouvelle liaison $(id, pk')$ et en
révoquant l'ancienne par un enregistrement de révocation signé.  La propriété
d'ajout exclusif implique que les révocations sont elles-mêmes permanentes et
publiquement vérifiables.

%--------------------------------------------------------------------
\section{Conclusion}
%--------------------------------------------------------------------

Le Socle Identitaire démontre qu'un substrat cryptographique minimal---un
schéma de signature EUF-AAMC et un annuaire en ajout exclusif avec finalité
quantifiée---est suffisant non seulement pour l'authentification, mais pour
permettre un ensemble complet de capacités système sans introduction d'ancres
de confiance institutionnelles.

Le modèle formel, la preuve de sécurité et l'analyse du coût d'attaque
économique établissent que les seules attaques significatives sur le système
exigent soit de violer des hypothèses d'intractabilité standard, soit de payer
un coût économique publiquement quantifiable.  Le système est déployé et
opérationnel depuis 2020 sur du matériel grand public, validant la faisabilité
pratique de la construction.

En dérivant le nommage, l'adressage par contenu, l'exécution déterministe, la
connectivité et l'interaction économique d'un unique primitif identitaire,
l'architecture réduit la complexité conceptuelle tout en préservant la
compatibilité avec les protocoles et normes existants.

%--------------------------------------------------------------------
\section*{Remerciements}
%--------------------------------------------------------------------

L'auteur remercie Oleg Khovayko pour les travaux fondateurs sur l'infrastructure
Emercoin et le cadre hack51 d'analyse du coût d'attaque~\cite{hack51}~;
Aleksej Sokolov (Alex014) pour la conception du format WORM~\cite{worm} et la
mise en \oe{}uvre de la chaîne d'outils PrivateNess (\texttt{privatenesstools},
\texttt{pyuheprng}) à partir des spécifications de l'auteur~; Brandon
\og~Synth~\fg~Smietana et Moses Narrow (\texttt{moe\_narrow} /
\texttt{0pcom}), responsable du packaging Linux Skywire, pour le transport en
réseau maillé Skywire et l'infrastructure de consensus Fiber sur lesquels NESS
est construit~; et la communauté PrivateNess Network, dont le travail collectif
depuis la fondation publique du projet sur Bitcointalk le 1\textsuperscript{er}
septembre 2020~\cite{bitcointalk} a rendu ce système déployé possible.  Neuf
années d'engagement de l'auteur dans l'infrastructure Emercoin constituent le
fondement opérationnel sur lequel le Socle Identitaire est bâti.

%--------------------------------------------------------------------
\begin{thebibliography}{99}

\bibitem{rompel1990}
  J.~Rompel,
  \og~One-way functions are necessary and sufficient for secure signatures,~\fg
  in \textit{Proc.\ 22\textsuperscript{e} Symposium annuel ACM sur la théorie
  de l'informatique (STOC)}, pp.~387--394, 1990.

\bibitem{haitner2010}
  I.~Haitner, T.~Holenstein, O.~Reingold, S.~Vadhan et H.~Wee,
  \og~Universal one-way hash functions via inaccessible entropy,~\fg
  in H.~Gilbert (éd.), \textit{EUROCRYPT 2010},
  LNCS vol.~6110, pp.~616--637, Springer, 2010.

\bibitem{naoryung1989}
  M.~Naor et M.~Yung,
  \og~Universal one-way hash functions and their cryptographic applications,~\fg
  in \textit{Proc.\ 21\textsuperscript{e} Symposium annuel ACM STOC},
  pp.~33--43, 1989.

\bibitem{katz2005}
  J.~Katz et C.-Y.~Koo,
  \og~On constructing universal one-way hash functions from arbitrary one-way
  functions,~\fg
  Cryptology ePrint Archive, Rapport 2005/328, 2005.
  \url{https://eprint.iacr.org/2005/328}

\bibitem{goldreich2001}
  O.~Goldreich,
  \textit{Foundations of Cryptography, Volume~1 : Basic Tools},
  Cambridge University Press, 2001.

\bibitem{eddsa}
  S.~Josefsson et I.~Liusvaara,
  \og~Edwards-Curve Digital Signature Algorithm (EdDSA),~\fg
  RFC~8032, IETF, janvier 2017.
  \url{https://www.rfc-editor.org/rfc/rfc8032}

\bibitem{fips204}
  National Institute of Standards and Technology,
  \og~Module-Lattice-Based Digital Signature Standard (ML-DSA),~\fg
  FIPS~204, août 2024.
  \url{https://doi.org/10.6028/NIST.FIPS.204}

\bibitem{uheprng}
  S.~Gibson,
  \og~Ultra-High Entropy Pseudo-Random Number Generator (UHEPRNG),~\fg
  Gibson Research Corporation.
  \url{https://www.grc.com/otg/uheprng.htm}

\bibitem{emercoin}
  Équipe de développement Emercoin,
  \og~Emercoin~: chaîne de blocs hybride avec minage fusionné,~\fg
  \url{https://emercoin.com}, 2013--2026.

\bibitem{hack51}
  O.~Khovayko,
  \og~Estimation du coût d'une attaque à 51\,\% sur Emercoin~: le coefficient
  hack51~\fg [en russe],
  Dzen / Yandex, 2021.
  \url{https://dzen.ru/a/YUotZkN9nAeElPEQ}

\bibitem{nuzzi2024}
  L.~Nuzzi et al.,
  \og~Breaking BFT~: quantifying the cost to attack Bitcoin and Ethereum,~\fg
  Coin Metrics, 2024.

\bibitem{skycoin}
  B.~Smietana (\og~Synth~\fg) et al.,
  \og~Skycoin~: Skywire et consensus Fiber,~\fg
  Équipe de développement Skycoin, 2014--2026.
  \url{https://www.skycoin.com}
  Packaging Linux~: M.~Narrow (\texttt{moe\_narrow} / \texttt{0pcom}).

\bibitem{worm}
  A.~Sokolov (Alex014) / NESS-Network,
  \og~WORM --- World Object MaRkup language,~\fg
  format conçu par A.~Sokolov~; déployé au sein de l'écosystème PrivateNess.
  \url{https://github.com/ness-network/worm}

\bibitem{onchain}
  PrivateNess Network,
  \og~PrivateNess.Network --- enregistrement d'identité canonique en
  chaîne,~\fg
  Emercoin NVS DPO, ancré dans le bloc~\#445\,669,
  1\textsuperscript{er}~novembre~2020, 22:16:45~UTC.  Opérationnel
  2020--2026.
  \url{https://explorer.emercoin.com/dpo/PrivateNESS.Network}

\bibitem{bitcointalk}
  PrivateNess Network,
  \og~PrivateNess Network (NESS) Privacy CX Fiber Coin,~\fg
  Bitcointalk Announcements (Altcoins), 1\textsuperscript{er}~septembre~2020.
  \url{https://bitcointalk.org/index.php?topic=5272741.0}

\bibitem{rfc5280}
  D.~Cooper et al.,
  \og~Internet X.509 Public Key Infrastructure Certificate and CRL Profile,~\fg
  IETF RFC~5280, mai 2008.

\bibitem{did2022}
  M.~Sporny et al.,
  \og~Decentralized Identifiers (DIDs) v1.0,~\fg
  Recommandation W3C, juillet 2022.
  \url{https://www.w3.org/TR/did-core/}

\bibitem{webauthn2021}
  D.~Balfanz et al.,
  \og~Web Authentication~: An API for Accessing Public Key Credentials
  Level~2,~\fg
  Recommandation W3C, avril 2021.
  \url{https://www.w3.org/TR/webauthn-2/}

\bibitem{eip4361}
  B.~Talsma et al.,
  \og~EIP-4361~: Sign-In with Ethereum,~\fg
  Proposition d'amélioration Ethereum, octobre 2021.
  \url{https://eips.ethereum.org/EIPS/eip-4361}

\bibitem{vc2019}
  M.~Sporny, D.~Longley et D.~Chadwick,
  \og~Verifiable Credentials Data Model 1.0,~\fg
  Recommandation W3C, novembre 2019.
  \url{https://www.w3.org/TR/vc-data-model/}

\bibitem{handshake}
  Auteurs Handshake,
  \og~Handshake~: Decentralized Certificate Authority and Naming,~\fg
  Article technique, 2018.
  \url{https://handshake.org/files/handshake.txt}

\bibitem{ipfs}
  J.~Benet,
  \og~IPFS---Content Addressed, Versioned, P2P File System,~\fg
  arXiv:1407.3561, 2014.

\end{thebibliography}

\end{document}
